\documentclass{situationist-preprint}

\title{The Reasons of Love}
\msid{EXI-16-a}
\msfamily{EXI-16}
\msdiscipline{Meaning, Death and the Absurd}
\mstype{Principal manuscript}
\mskeywords{love; reasons; fittingness; substitution; trading up; Kolodny; Frankfurt; Velleman; relationships}

\begin{document}
\maketitle

\begin{abstract}
Is love rational? If love is a response to the beloved's qualities, it seems that it should transfer to anyone with the same qualities and shift to anyone with better ones, which is not how love works and not how lovers want to be loved. If love is not a response to reasons at all, as Frankfurt holds, it cannot be criticized as unfitting, which it can. We argue that love responds to reasons, but that these reasons do not make love transferable to qualitatively similar persons. The key move is to distinguish two kinds of reasons for love: qualitative reasons, which are features of the beloved that anyone with those features would share, and relational reasons, which consist in the history of a relationship between the lover and the beloved. We show that love is answerable to both, that qualitative reasons alone would make love transferable, and that relational reasons, which no qualitatively similar person shares, block the transfer. Qualitative reasons make the beginning of love intelligible; relational reasons sustain it and make it answerable to how the relationship goes. The result develops Kolodny's relationship view while taking the person, not the relationship, to be the object of love.
\end{abstract}

\section{The problem}

The oldest answer to the question why we love is that we love what is good or beautiful in the beloved. In the speech of Diotima, the lover who has loved one beautiful body comes to see that the beauty of any body is akin to that of any other, and so to love all beautiful bodies, and then beauty itself \citep[210a--212a]{PlatoSymp}. The ascent is a reductio of the view for those who think love is directed at a particular person: if the reasons for love are qualities, the person is loved only as an instance of them. Two problems follow. The \emph{substitution problem} is that a qualitative duplicate of the beloved would be an equally fitting object of love, so that love could be transferred to the duplicate without loss. The \emph{trading-up problem} is that a person with more of the relevant qualities would be a more fitting object, so that love ought to shift to her.

Responses fall into three groups. \citet{Frankfurt2004} denies that love is a response to reasons: love is a configuration of the will, a concern for the beloved's good for its own sake, which is not grounded in any judgment of the beloved's value; it creates reasons rather than answering to them. \citet{Velleman1999} holds that love responds to the beloved's value as a person, the rational nature that all persons share, and that its selectivity is explained by the contingent fit between the beloved's empirical self and the lover's capacity to see that value. \citet{Kolodny2003} holds that love is valuing a relationship: the reason for loving someone is the history of the relationship one has with her, and love is the disposition to value both the relationship and the person as one's partner in it. Among property views, \citet{Keller2000} defends the claim that the beloved's properties are reasons for love and argues that the view is less vulnerable to the substitution problem than it appears; \citet{Jollimore2011} treats love as a mode of attention to the beloved's valuable features. \citet{Delaney1996} identifies what lovers want: to be loved for properties they take to matter, and to be loved in a way that does not pass to anyone who has them. \citet{Setiya2014} objects to relationship views that one can love a person with whom one has no relationship, and locates the reason in the beloved's humanity.

We argue for the conclusion that love responds to reasons, but that these reasons do not make love transferable to qualitatively similar persons.

\section{Framework}

\begin{definition}[Reasons for love]\label{d:reason}
A \emph{reason for love} of $x$ by $L$ is a consideration that counts in favor of $L$'s loving $x$, in the sense that it makes the love fitting and can be cited in its justification, and whose absence or loss counts against it.
\end{definition}

\begin{definition}[Qualitative and relational reasons]\label{d:kinds}
A reason for $L$'s love of $x$ is \emph{qualitative} if it is a feature of $x$ that does not involve $L$, such as kindness, wit or courage. It is \emph{relational} if it consists in a relationship between $L$ and $x$ constituted by a shared history: past interactions, mutual concern, commitments made and kept.
\end{definition}

\begin{definition}[Qualitative similarity]\label{d:sim}
A person $y$ is \emph{qualitatively similar} to $x$ relative to $L$ if $y \neq x$ and $y$ has every qualitative feature of $x$ that is a reason for $L$'s love of $x$, to at least the same degree.
\end{definition}

\begin{definition}[Transferability]\label{d:trans}
$L$'s love of $x$ is \emph{transferable} if, for every person $y$ qualitatively similar to $x$ relative to $L$, the reasons for $L$'s love of $x$ would be at least as strong reasons for $L$'s loving $y$ instead, so that $L$ could love $y$ in place of $x$ with no loss in what her love answers to.
\end{definition}

\begin{principle}[Answerability]\label{p:ans}
An attitude responds to reasons if it is subject to rational criticism for being unfitting to its object, and if its persistence or withdrawal can be made intelligible by changes in the considerations that bear on it.
\end{principle}

\begin{principle}[Generality]\label{p:gen}
If a consideration that is a feature of $x$ is a reason for an attitude towards $x$, the same feature, had to the same degree by $y$, is an equally strong reason for the same attitude towards $y$.
\end{principle}

\begin{principle}[Particularity of history]\label{p:hist}
A relationship constituted by a shared history between $L$ and $x$ is not had by $L$ and any $y \neq x$: a numerically distinct person, however qualitatively similar, did not take part in the events that constitute it.
\end{principle}

Principle~\ref{p:ans} is the usual mark of reason-responsiveness: beliefs, fears and resentments are criticizable as unfitting and change intelligibly with the evidence. Principle~\ref{p:gen} is the generality of reasons that generates both problems: it applies to features of the object. Principle~\ref{p:hist} is a truism about histories, which consist of particular events with particular participants; a qualitatively similar history between $L$ and $y$ would be another history.

\section{The argument}

\begin{lemma}\label{l:resp}
Love responds to reasons.
\end{lemma}

\begin{proof}
We show that love meets both conditions of Principle~\ref{p:ans}. First, love is criticizable as unfitting. A person who loves a partner who has systematically deceived and exploited her, and who continues to love him after learning this, is not merely unlucky; her friends may rightly say that he does not deserve her love and that her love is misplaced, and she may come to say so herself. A person who loves someone she has never met, on the strength of an image she has constructed, loves something that does not answer to the person. These criticisms presuppose that some considerations count for and against love. Second, the course of love is intelligible in terms of such considerations. Love that ends because the beloved has become cruel, or because the relationship has been destroyed by betrayal, ends for a reason, and the lover can cite it; love that persists through illness and change can be explained by the history it rests on. If love were wholly independent of reasons, as \citet{Frankfurt2004} holds, these criticisms would be confused and these explanations unavailable.
\end{proof}

\begin{lemma}\label{l:qual}
If all reasons for $L$'s love of $x$ are qualitative, the love is transferable.
\end{lemma}

\begin{proof}
Let $y$ be qualitatively similar to $x$ relative to $L$. By Definition~\ref{d:sim}, $y$ has every qualitative feature that is a reason for $L$'s love of $x$, to at least the same degree. By Principle~\ref{p:gen}, each of these features is at least as strong a reason for $L$'s loving $y$. If there are no other reasons, the reasons for loving $x$ are at least as strong reasons for loving $y$ instead, and by Definition~\ref{d:trans} the love is transferable. If $y$ has some features to a greater degree, the same argument yields trading up.
\end{proof}

\begin{lemma}\label{l:rel}
If the reasons for $L$'s love of $x$ include a relational reason, the love is not transferable.
\end{lemma}

\begin{proof}
Let $R$ be the relationship that is a relational reason for $L$'s love of $x$, and let $y$ be qualitatively similar to $x$. By Principle~\ref{p:hist}, $L$ does not have $R$ with $y$. Qualitative similarity concerns only features of $x$ that do not involve $L$ (Definitions~\ref{d:kinds} and~\ref{d:sim}), so it does not supply $y$ with $R$, and Principle~\ref{p:gen} does not apply, since $R$ is not a feature had by $x$ alone. Loving $y$ in place of $x$ would therefore lose a reason that $L$'s love of $x$ answers to. By Definition~\ref{d:trans}, the love is not transferable.
\end{proof}

\begin{lemma}\label{l:sust}
Every established love answers to relational reasons.
\end{lemma}

\begin{proof}
A love is established when it has persisted through a history of interaction with the beloved. Such a love is answerable to that history: it can be criticized as unfitting when the history is one of exploitation, and its ending is made intelligible by betrayal or estrangement, which are events in the relationship rather than changes in the beloved's qualities (Lemma~\ref{l:resp}). Conversely, the lover can cite the history as what she loves her beloved for: the years of mutual care, the trust built and kept. By Definition~\ref{d:reason}, the history is a reason for the love, and by Definition~\ref{d:kinds} it is relational.
\end{proof}

\begin{theorem}\label{t:main}
Love responds to reasons, but these reasons do not make love transferable to qualitatively similar persons.
\end{theorem}

\begin{proof}
By Lemma~\ref{l:resp}, love responds to reasons. By Lemma~\ref{l:sust}, every established love answers to relational reasons, and by Lemma~\ref{l:rel}, a love that answers to a relational reason is not transferable. Love that is not yet established is not yet a love that could be transferred; what precedes it is attraction, interest or the beginning of a relationship, and these, as Lemma~\ref{l:qual} predicts, are directed at qualities and could have been directed at anyone who had them.
\end{proof}

The theorem assigns the two kinds of reason different roles. Qualitative reasons make the beginning of love intelligible: we come to love people because they are kind, or funny, or brave, and to ask why someone began to love another is to ask for such features. Relational reasons sustain love and make it answerable to how the relationship goes. This explains Delaney's datum. Lovers want to be loved for properties they care about, which qualitative reasons provide, and they want the love not to pass to whoever has those properties, which relational reasons secure \citep{Delaney1996}. It also explains why trading up is not a requirement of reason: a person with better qualities offers better qualitative reasons, but no history, and the reasons of an established love are not simply outweighed by a surplus of qualities, since what the love answers to is not on offer elsewhere.

The result fits the account of partiality in ETH-19, which argues that agents may give priority to those close to them without justifying the priority by impartial reasons. The reasons of love are reasons about this person and this relationship, and they are not certified by any consideration that would apply equally to a qualitatively similar stranger.

\section{Objections and replies}

\begin{objection}[Wrong object]\label{o:obj}
On the relationship view, the lover values the relationship, and the beloved only as a participant in it. But we love people, not relationships, and a love whose object is the relationship is in danger of being self-regarding \citep{Setiya2014}.
\end{objection}

\begin{reply}
Theorem~\ref{t:main} concerns the reasons of love, not its object. The object of love is the beloved; the relationship is among the considerations in virtue of which loving her is fitting, as a friend's past kindness is a reason for gratitude whose object is the friend. We therefore take from \citet{Kolodny2003} the claim that relationships supply reasons for love, without the claim that love is a form of valuing the relationship. Setiya's case of love without a relationship, such as love at first sight, is on our account the beginning of love, answerable to qualitative reasons and transferable in the sense of Lemma~\ref{l:qual}; this is the correct verdict, since nothing would be lost if the first sight had been of a qualitative duplicate.
\end{reply}

\begin{objection}[Duplication]\label{o:dup}
Suppose the beloved is duplicated, molecule for molecule, together with her memories of the relationship. The duplicate remembers the shared history as well as the original does, so relational reasons do not discriminate between them.
\end{objection}

\begin{reply}
Remembering a history is a qualitative feature of the rememberer; having taken part in it is not. If the duplicate is the beloved, by whatever account of personal identity is correct (PID-01 proposes psychological continuity with bodily continuity as a tie-breaker in cases of branching), the love is not transferred but continued. If the duplicate is not the beloved, she has apparent memories of events in which she did not take part, and she does not share the history that the love answers to. Either way, the case is not one in which love passes to a qualitatively similar other person with no loss in its reasons.
\end{reply}

\begin{objection}[Arbitrariness]\label{o:arb}
If an established love rests on a history, and the history began with qualities that many people share, the love rests in the end on chance: on whom one happened to meet. Its reasons are therefore no better than Frankfurt's brute volitional facts.
\end{objection}

\begin{reply}
That a reason arose by chance does not make it less of a reason. A promise made to the person one happened to meet binds as fully as any other, and the obligations of a parent to a child are not weakened by the contingency of which child was born. The history of a relationship, once it exists, is a fact about the lover and the beloved to which love can be more or less fitting, and that is all Lemma~\ref{l:resp} requires.
\end{reply}

\begin{objection}[Universal value]\label{o:vell}
Following \citet{Velleman1999}, the reason for love is the beloved's value as a person, which is the same in every person. This reason is not qualitative in the problematic sense and does not license substitution, since it does not rank persons. Relational reasons are unnecessary.
\end{objection}

\begin{reply}
The beloved's value as a person may well be a reason for love, but it is equally a reason for loving every person, and so cannot explain why love is answerable to the course of a particular relationship. On Velleman's view, a love that ends because the relationship has been betrayed does not end for a reason, since the beloved's rational nature is unchanged. Lemma~\ref{l:resp} shows that it does end for one. Value as a person is at most a background condition that makes love of persons fitting at all.
\end{reply}

\section{What the result does not establish}

The argument does not show that love is ever required by reasons, or that a person who fails to love someone with good qualities or a long shared history is irrational; reasons for love may be reasons that make it fitting without making it mandatory. It does not settle whether love is an emotion, a volitional structure or a form of valuing; the theorem requires only that love be answerable to considerations in the sense of Principle~\ref{p:ans}. Finally, it does not show that qualitative reasons are irrelevant to established love: a profound change in the beloved can make the end of a love intelligible, and Lemma~\ref{l:rel} claims only that the relational reasons that remain prevent the love from passing to someone else.

\section{Conclusion}

Love is rational in the sense that it responds to reasons: it can be unfitting, and its course is intelligible in terms of considerations the lover can cite. Its reasons are of two kinds. Qualitative reasons make its beginning intelligible and, taken alone, would make it transferable to anyone with the same qualities. Relational reasons, which consist in a shared history that no other person has with the lover, sustain every established love, and they make it non-transferable. Love responds to reasons, but these reasons do not make love transferable to qualitatively similar persons.

\begin{thebibliography}{9}\small
\bibitem[Delaney(1996)]{Delaney1996} Delaney, N. (1996). Romantic love and loving commitment: Articulating a modern ideal. \emph{American Philosophical Quarterly}, 33(4), 339--356.
\bibitem[Frankfurt(2004)]{Frankfurt2004} Frankfurt, H.~G. (2004). \emph{The Reasons of Love}. Princeton University Press.
\bibitem[Jollimore(2011)]{Jollimore2011} Jollimore, T. (2011). \emph{Love's Vision}. Princeton University Press.
\bibitem[Keller(2000)]{Keller2000} Keller, S. (2000). How do I love thee? Let me count the properties. \emph{American Philosophical Quarterly}, 37(2), 163--173.
\bibitem[Kolodny(2003)]{Kolodny2003} Kolodny, N. (2003). Love as valuing a relationship. \emph{The Philosophical Review}, 112(2), 135--189.
\bibitem[Plato(Symp.)]{PlatoSymp} Plato. \emph{Symposium}. Cited by Stephanus pages.
\bibitem[Setiya(2014)]{Setiya2014} Setiya, K. (2014). Love and the value of a life. \emph{The Philosophical Review}, 123(3), 251--280.
\bibitem[Velleman(1999)]{Velleman1999} Velleman, J.~D. (1999). Love as a moral emotion. \emph{Ethics}, 109(2), 338--374.
\end{thebibliography}

\end{document}
